\centerline{\bf \Huge Тригонометрия с ЕГЭ}


\centerline{\bf \Huge 2025}

\problem{\textbf{1}} а) Решите уравнение $2 \sin (-x)+2 \sqrt{3} \sin x-4 \cos ^2 x=\sqrt{3}-4$.

б) Найдите все корни этого уравнения, принадлежащие отрезку $\left[-5 \pi ;-\dfrac{7 \pi}{2}\right]$.

\problem{\textbf{2}} a) Решите уравнение $2-2 \cos (\pi+2 x)-\sqrt{8} \cos x=\sqrt{6}-\sqrt{12} \cos x$.

б) Укажите корни этого уравнения, принадлежащие отрезку $\left[\dfrac{\pi}{2} ; 2 \pi\right]$.

\centerline{\bf \Huge 2024}

\problem{\textbf{3}} а) Решите уравнение $2 \cos x+\sin ^2 x=2 \cos ^3 x$.

б) Определите, какие из его корней принадлежат отрезку $\left[-\dfrac{9 \pi}{2} ;-3 \pi\right]$.

\problem{\textbf{4}} a) Решите уравнение $\cos ^2(\pi-x)-\sin \left(\dfrac{3 \pi}{2}+x\right)=0$.

б) Найдите все корни этого уравнения, принадлежащие отрезку $\left[-2 \pi ;-\dfrac{\pi}{2}\right]$.

\centerline{\bf \Huge 2023}

\problem{\textbf{5}} а) Решите уравнение $2 \cos ^3 x=\sqrt{3} \sin ^2 x+2 \cos x$.

б) Найдите все корни этого уравнения, принадлежащие отрезку $\left[-3 \pi ;-\dfrac{3 \pi}{2}\right]$.

\problem{\textbf{6}} a) Решите уравнение $\sin x \cdot \cos 2 x+\sin x=\sqrt{3} \cos ^2 x$.

б) Найдите все корни этого уравнения, принадлежащие отрезку $\left[\dfrac{7 \pi}{2} ; \frac{9 \pi}{2}\right]$.

\centerline{\bf \Huge 2022}

\problem{\textbf{7}} а) Решите уравнение: $81^{\cos x}-12 \cdot 9^{\cos x}+27=0$.

б) Определите, какие из его корней принадлежат отрезку $\left[-4 \pi ;-\dfrac{5 \pi}{2}\right]$.

\problem{\textbf{8}} a) Решите уравнение $2 \sin 2 x+2 \sin (-x)-2 \cos (-x)+1=0$.

б) Укажите корни этого уравнения, принадлежащие отрезку $\left[\dfrac{5 \pi}{2} ; 4 \pi\right]$.



